\chapter{Survival laws and the atlas of obstructions}\label{ch:atlas}

\Cref{ch:topology} analysed one explicit cycle. The census of \cref{obs:census-H2} shows that this cycle explains only part of the $H_2$ that SELECT-II produces. The structure behind the anchor is in fact general: for \emph{any} cycle of the permutohedron the survival law has the same form, and the family of laws that occur is itself an object of study. The theorem of \cref{sec:atlas-law} is general; the catalogue of \cref{sec:atlas-census} is a finite census with an explicit prediction record.

\section{The survival law of an arbitrary cycle}\label{sec:atlas-law}
Fix $k$ and write $A(m)=\mathfrak F_\lambda(m)\subseteq 2^{[k]}\setminus\{\emptyset\}$ for the family of feasible prefix sets of a content at source $m$ (\cref{lem:F}). For a $j$-cycle $z$ of $P_k$ put $P_z=\bigcup_{F\in\operatorname{supp}z}\Pref(F)$ and $K_z(\lambda)=\lcm_{I\in P_z}D(I)$.
@@def:filling-family@@
Filling is monotone: enlarging $S$ enlarges the set of available $(j{+}1)$-faces, hence the image of $\partial_{j+1}$. So the filling sets form an up-set, determined by the antichain $\mathfrak M(z)$.
@@thm:AT1@@
Nothing arithmetic enters beyond the monotonicity of $A(m)$ in $m$ under divisibility. The content of the theorem is that \emph{every} survival question reduces to a monotone Boolean function of at most $2^k$ prefix sets, whose minimal true points $\mathfrak M(z)$ form a finite combinatorial invariant of $(P_k,z)$ alone. Arithmetic enters only through which sets lie in $A(m)$ (\cref{prop:valuation}), and \cref{prop:CRT} says how freely these can be prescribed. For the anchor, $\mathfrak M(z)=\{B_0,B_1,B_2,B_3\}$ (\cref{thm:exact-survival}); \cref{thm:hole-survival} is the weaker statement that follows from the fact that each $B_i$ contains a gate.

\paragraph{Computing $\mathfrak M(z)$.} The set of non-filling families is a down-set whose maximal elements $N_1,\dots,N_t$ determine $\mathfrak M(z)$: a set fills iff it meets the complement of every $N_i$, so $\mathfrak M(z)$ is the set of minimal transversals of $\{\overline{N_i}\}$. The computation alternates Berge dualisation with greedy extension of a non-filling transversal to a maximal non-filling set, and stops when every minimal transversal tests as filling; it is exact, with no enumeration cut-off.

\section{The atlas}\label{sec:atlas-census}
@@obs:atlas+finite-exact@@
@@obs:atlas-pred+bounded@@

\section{Scope of the atlas}
Every atlas statement is relative to six-label contents, the field $\F$ and the SELECT-II family $\mathfrak F_\lambda(m)=\{I:D(I)\mid m\}$.
@@prop:atlas-scope@@
Three limits of the catalogue should be kept in view. The number of law classes is not shown to be finite for six labels; the observed rates of new classes on successive held-out rounds ($18/211$, $13/171$, $3/292$) decrease, but no bound is proved. Classes are defined from low-weight representatives of homology classes, and different representatives of one class may have different laws, so whether the $32$ classes merge when defined from the homology class itself is not settled. And although \cref{prop:gate-nec} shows that the large families of the anchor class can be arithmetically realised or avoided, the corresponding realisability question for the other classes has not been tested.
